\documentclass[11pt,a4paper,oneside]{report}

% -----------------------------------------------------------------------------
% Typography and page geometry
% -----------------------------------------------------------------------------
\usepackage[T1]{fontenc}
\usepackage[utf8]{inputenc}
% Use the installed Times fonts when NewTX is unavailable.
\newif\ifnewtxavailable
\IfFileExists{newtxtext.sty}{\IfFileExists{newtxmath.sty}{\newtxavailabletrue}{}}{}
\ifnewtxavailable
  \usepackage{newtxtext}
\else
  \usepackage{mathptmx}
\fi
\usepackage[scaled=0.94]{helvet}
\usepackage[a4paper,left=2.45cm,right=2.45cm,top=2.45cm,bottom=2.55cm,headheight=16pt,headsep=0.55cm,footskip=1.15cm]{geometry}
\usepackage{microtype}
\usepackage{ragged2e}
\usepackage{needspace}

% -----------------------------------------------------------------------------
% Mathematics, tables, graphics, and layout
% -----------------------------------------------------------------------------
\usepackage{amsmath,mathtools}
\ifnewtxavailable
  \usepackage{newtxmath}
\else
  \usepackage{amssymb}
\fi
\usepackage{booktabs,tabularx,array,longtable,multirow}
\usepackage{enumitem}
\usepackage[table]{xcolor}
\usepackage[skins,breakable]{tcolorbox}
\usepackage{graphicx}
\graphicspath{{cv_assets/}{./}}
% Keep the notes buildable when the separately supplied figure assets are absent.
% Use the same search order as graphicspath and restore each image automatically.
\newcommand{\lectureimage}[3][0.36\textheight]{%
  \IfFileExists{cv_assets/#3}{%
    \includegraphics[width=#2,height=#1,keepaspectratio]{cv_assets/#3}%
  }{%
    \IfFileExists{#3}{%
      \includegraphics[width=#2,height=#1,keepaspectratio]{#3}%
    }{%
      \PackageWarning{ComputerVision}{Missing figure: #3}%
      \parbox[c][2cm][c]{#2}{\centering\small\sffamily
        Figure unavailable\par\medskip
        \texttt{\detokenize{#3}}}%
    }%
  }%
}
\usepackage{tikz}
\usetikzlibrary{arrows.meta,positioning,fit,calc,shapes.geometric,backgrounds,matrix,decorations.pathreplacing}
\usepackage{pgfplots}
\pgfplotsset{compat=1.18}
\usepackage{caption}
\usepackage{fancyhdr}
\usepackage{titlesec}
\usepackage{etoolbox}
\makeatletter
\patchcmd{\chapter}{\if@openright\cleardoublepage\else\clearpage\fi}
  {\FloatBarrier\par\Needspace{16\baselineskip}}{}{}
\makeatother
\usepackage{tocloft}
\usepackage{hyperref}
\usepackage{bookmark}
\usepackage{multicol}
\usepackage{float}
\usepackage{placeins}
\renewcommand{\topfraction}{0.88}
\renewcommand{\bottomfraction}{0.75}
\renewcommand{\textfraction}{0.08}
\renewcommand{\floatpagefraction}{0.70}
\tikzset{cvnode/.style={draw=LightGrey,fill=PaleGrey,rounded corners=1mm,
  minimum height=10mm,text width=29mm,align=center,font=\small\sffamily},
  cvarrow/.style={-{Latex[length=2mm]},line width=0.7pt,draw=SapienzaRed}}

% Keep figures and tables in reading order with their surrounding discussion.
\setlength{\intextsep}{10pt plus 2pt minus 2pt}

% -----------------------------------------------------------------------------
% Colour system
% -----------------------------------------------------------------------------
\definecolor{SapienzaRed}{HTML}{7A1737}
\definecolor{SapienzaDark}{HTML}{5B1028}
\definecolor{DeepGreen}{HTML}{3F6B55}
\definecolor{MutedBlue}{HTML}{416A86}
\definecolor{MutedAmber}{HTML}{A66C20}
\definecolor{Ink}{HTML}{202020}
\definecolor{MidGrey}{HTML}{555555}
\definecolor{LightGrey}{HTML}{E5E5E5}
\definecolor{PaleGrey}{HTML}{F7F7F7}
\definecolor{PaleRed}{HTML}{F7EFF2}
\definecolor{PaleGreen}{HTML}{F1F6F3}
\definecolor{PaleBlue}{HTML}{EFF4F7}
\definecolor{PaleAmber}{HTML}{FBF5EA}
\definecolor{TableTint}{HTML}{F4E9ED}

\hypersetup{
  colorlinks=true,
  linkcolor=SapienzaRed,
  urlcolor=SapienzaRed,
  citecolor=SapienzaRed,
  pdftitle={Computer Vision - Lecture Notes},
  pdfauthor={Matt Clifford Trinidad Ramos},
  pdfsubject={Computer Vision lecture notes: introduction, image acquisition, color, and texture},
  pdfkeywords={computer vision, image acquisition, sampling, quantization, color, white balance, texture, LBP}
}

\setlength{\parindent}{0pt}
\setlength{\parskip}{0.46em}
\raggedbottom
\setlist[itemize]{topsep=0.35em,itemsep=0.22em,parsep=0pt,leftmargin=1.55em}
\setlist[enumerate]{topsep=0.35em,itemsep=0.22em,parsep=0pt,leftmargin=1.75em}
\setcounter{tocdepth}{1}

% -----------------------------------------------------------------------------
% Chapter and section styling
% -----------------------------------------------------------------------------
\titleformat{\chapter}[display]
  {\normalfont\sffamily\bfseries\color{Ink}\raggedright\hyphenpenalty=10000\exhyphenpenalty=10000}
  {\fontsize{10.5}{12}\selectfont\color{SapienzaRed}\MakeUppercase{Chapter \thechapter}}
  {0.45em}
  {\fontsize{23}{27}\selectfont}
  [\vspace{0.68em}{\color{SapienzaRed}\titlerule[0.8pt]}]
\titlespacing*{\chapter}{0pt}{20pt}{19pt}

\titleformat{\section}
  {\normalfont\sffamily\Large\bfseries\color{Ink}}
  {\color{SapienzaRed}\thesection}{0.70em}{}
\titlespacing*{\section}{0pt}{1.35em}{0.48em}

\titleformat{\subsection}
  {\normalfont\sffamily\large\bfseries\color{Ink}}
  {\color{SapienzaRed}\thesubsection}{0.65em}{}
\titlespacing*{\subsection}{0pt}{0.95em}{0.34em}

% -----------------------------------------------------------------------------
% Table of contents styling
% -----------------------------------------------------------------------------
\renewcommand{\cfttoctitlefont}{\sffamily\Huge\bfseries\color{Ink}}
\renewcommand{\cftchapfont}{\sffamily\bfseries\color{Ink}}
\renewcommand{\cftchappagefont}{\sffamily\bfseries\color{SapienzaRed}}
\renewcommand{\cftsecfont}{\normalsize\sffamily\color{Ink}}
\renewcommand{\cftsecpagefont}{\normalsize\sffamily\color{MidGrey}}
\renewcommand{\cftsubsecfont}{\small\sffamily\color{MidGrey}}
\renewcommand{\cftsubsecpagefont}{\small\sffamily\color{MidGrey}}
\setlength{\cftbeforechapskip}{0.45em}
\setlength{\cftbeforesecskip}{0.04em}

% -----------------------------------------------------------------------------
% Reusable academic boxes
% -----------------------------------------------------------------------------
\newtcolorbox{definitionbox}[1]{
  enhanced,breakable,frame hidden,
  colback=PaleRed,
  borderline west={1.8pt}{0pt}{SapienzaRed},
  title={#1},
  coltitle=SapienzaRed,
  fonttitle=\sffamily\bfseries,
  left=3.2mm,right=3.2mm,top=2.0mm,bottom=2.0mm,
  before skip=0.8em,after skip=0.8em
}

\newtcolorbox{keyconcept}[1][Key concept]{
  enhanced,breakable,frame hidden,
  colback=PaleGreen,
  borderline west={1.8pt}{0pt}{DeepGreen},
  title={#1},
  coltitle=DeepGreen,
  fonttitle=\sffamily\bfseries,
  left=3.2mm,right=3.2mm,top=1.6mm,bottom=1.6mm,
  before skip=0.7em,after skip=0.7em
}

\newtcolorbox{remarkbox}[1]{
  enhanced,breakable,frame hidden,
  colback=PaleGrey,
  borderline west={1.4pt}{0pt}{MidGrey},
  title={#1},
  coltitle=MidGrey,
  fonttitle=\sffamily\bfseries,
  left=3.2mm,right=3.2mm,top=2.0mm,bottom=2.0mm,
  before skip=0.8em,after skip=0.8em
}

\newtcolorbox{examplebox}[1]{
  enhanced,breakable,frame hidden,
  colback=PaleGrey,
  borderline west={1.6pt}{0pt}{SapienzaRed},
  title={#1},
  coltitle=SapienzaRed,
  fonttitle=\sffamily\bfseries,
  left=3.2mm,right=3.2mm,top=2.0mm,bottom=2.0mm,
  before skip=0.75em,after skip=0.75em
}

\newtcolorbox{chaptersummary}{
  enhanced,breakable,frame hidden,
  colback=PaleGrey,
  borderline west={1.6pt}{0pt}{SapienzaRed},
  left=3.0mm,right=3.0mm,top=1.15mm,bottom=1.15mm,
  before skip=0.65em,after skip=0.10em,
  before upper={\textcolor{SapienzaRed}{\sffamily\bfseries Chapter summary.}\enspace}
}


\newtcolorbox{warningbox}[1][Common confusion]{
  enhanced,breakable,frame hidden,
  colback=PaleAmber,
  borderline west={1.7pt}{0pt}{MutedAmber},
  title={#1}, coltitle=MutedAmber,
  fonttitle=\sffamily\bfseries,
  left=3.2mm,right=3.2mm,top=1.8mm,bottom=1.8mm,
  before skip=0.75em,after skip=0.75em
}

\newtcolorbox{metricbox}[1][Key relation]{
  enhanced,breakable,frame hidden,
  colback=PaleBlue,
  borderline west={1.7pt}{0pt}{MutedBlue},
  title={#1}, coltitle=MutedBlue,
  fonttitle=\sffamily\bfseries,
  left=3.2mm,right=3.2mm,top=1.8mm,bottom=1.8mm,
  before skip=0.75em,after skip=0.75em
}

\newcommand{\lecturedivider}[4]{%
  \clearpage
  \thispagestyle{empty}
  \phantomsection
  \ifnum#1=2\relax\addtocontents{toc}{\protect\newpage}\fi
  \addcontentsline{toc}{part}{Lecture #1: #2}
  \vspace*{0.22\textheight}
  {\sffamily\small\bfseries\color{SapienzaRed}LECTURE #1\par}
  \vspace{0.45em}
  {\sffamily\fontsize{28}{33}\selectfont\bfseries\color{Ink}#2\par}
  \vspace{0.8em}
  {\color{SapienzaRed}\rule{0.68\textwidth}{1pt}\par}
  \vspace{1.0em}
  {\large\color{MidGrey}#3\par}
  \vspace{2.2em}
  {\sffamily\color{MidGrey}#4\par}
  \vspace{3em}
  {\color{SapienzaRed}\rule{2cm}{3pt}\par}
  \vfill
  {\color{MidGrey}Computer Vision --- Marco Raoul Marini\par}
  \vfill
  {\small\color{MidGrey}Notes prepared by Matt Clifford Trinidad Ramos\par}
  \clearpage
}

\newcommand{\mathlegend}[1]{%
  \par\vspace{0.2em}{\small\itshape where}\par
  \begin{itemize}[leftmargin=1.7em,topsep=0.15em,itemsep=0.12em]
  #1
  \end{itemize}
}

\newcommand{\reviewgroup}[1]{%
  \par\vspace{0.7em}%
  {\sffamily\bfseries\color{SapienzaRed}#1}\par\vspace{0.10em}%
}

% -----------------------------------------------------------------------------
% Headers, footers, tables, captions
% -----------------------------------------------------------------------------
\pagestyle{fancy}
\fancyhf{}
\fancyhead[L]{\small\sffamily\color{MidGrey}\nouppercase{\leftmark}}
\fancyhead[R]{\small\sffamily\color{MidGrey}Computer Vision}
\fancyfoot[L]{\footnotesize\sffamily\color{MidGrey}Computer Vision}
\fancyfoot[C]{\footnotesize\sffamily\color{SapienzaRed}\thepage}
\fancyfoot[R]{\footnotesize\sffamily\color{MidGrey}Matt Clifford Trinidad Ramos}
\renewcommand{\headrulewidth}{0.3pt}
\renewcommand{\footrulewidth}{0pt}
\makeatletter
\renewcommand{\chaptermark}[1]{\markboth{Chapter \thechapter}{}}
\renewcommand{\sectionmark}[1]{\markright{#1}}
\makeatother

\captionsetup{font=small,labelfont={bf,sf,color=SapienzaRed},textfont={small},labelsep=period}
\captionsetup[table]{position=top,skip=5pt}
\renewcommand{\arraystretch}{1.25}
\newcolumntype{Y}{>{\RaggedRight\arraybackslash}X}

% -----------------------------------------------------------------------------
% Document
% -----------------------------------------------------------------------------
\begin{document}

\hypersetup{pageanchor=false}
\begin{titlepage}
  \thispagestyle{empty}
  \centering
  \vspace*{0.85cm}

  \IfFileExists{sapienza_logo.png}{%
    \includegraphics[width=0.50\textwidth,keepaspectratio]{sapienza_logo.png}\par
    \vspace{0.95cm}
  }{%
    {\sffamily\small\bfseries\color{SapienzaRed}SAPIENZA UNIVERSITY OF ROME\par}
    \vspace{1.15cm}
  }

  {\color{SapienzaRed}\rule{0.84\textwidth}{1.0pt}}\par
  \vspace{0.75cm}

  {\sffamily\fontsize{29}{33}\selectfont\bfseries\color{Ink}\MakeUppercase{Computer Vision}\par}
  \vspace{0.26cm}
  {\sffamily\Large\color{MidGrey}Lecture Notes\par}

  \vspace{0.75cm}
  {\color{SapienzaRed}\rule{0.84\textwidth}{1.0pt}}\par

  \vfill

  {\sffamily\footnotesize\bfseries\color{MidGrey}\MakeUppercase{Lecturer}\par}
  \vspace{0.14cm}
  {\large Marco Raoul Marini\par}

  \vspace{0.85cm}
  {\small Department of Computer Science\par}
  {\small Sapienza University of Rome\par}

  \vspace{1.35cm}
  {\sffamily\footnotesize\bfseries\color{MidGrey}\MakeUppercase{Notes prepared by}\par}
  \vspace{0.14cm}
  {\large\bfseries Matt Clifford Trinidad Ramos\par}

  \vfill
  {\color{LightGrey}\rule{0.84\textwidth}{0.5pt}}\par
\end{titlepage}
\hypersetup{pageanchor=true}

\pagenumbering{roman}
\begingroup
\setlength{\parskip}{0pt}
\normalsize
\markboth{Contents}{Contents}\tableofcontents
\endgroup
\clearpage
\pagenumbering{arabic}

\lecturedivider{1}{Introduction: Scope, Challenges, and Applications}{From visual measurements to descriptions of the world.}{Scope $\cdot$ Visual cues $\cdot$ Geometry $\cdot$ Recognition}

% =============================================================================
\chapter{Computer Vision: Purpose and Scope}
% =============================================================================

Computer vision aims to make computational systems extract properties of the three-dimensional world from one or more images. The desired result may be a physical measurement, an object or scene category, a location, a spatial relation, or another description that is useful for a task.

\begin{figure}[!htbp]
\centering
\fcolorbox{LightGrey}{white}{\lectureimage{0.88\textwidth}{computer_vision_eye.jpg}}
\caption{A visual metaphor for the central goal of computer vision: enabling computational systems to interpret visual input.}
\label{fig:making-computers-see}
\end{figure}


\begin{figure}[!htbp]
\centering
\begin{tikzpicture}[node distance=5mm]
\node[cvnode] (input) {Image / video};
\node[cvnode,right=of input] (sense) {Sensing device};
\node[cvnode,right=of sense] (interpret) {Interpreting device};
\node[cvnode,right=of interpret] (scene) {Interpretation\\of the scene};
\draw[cvarrow] (input)--(sense);
\draw[cvarrow] (sense)--(interpret);
\draw[cvarrow] (interpret)--(scene);
\end{tikzpicture}
\caption{The source's computer-vision pipeline: visual input becomes an interpretation of the world through sensing and interpretation.}
\label{fig:vision-pipeline}
\end{figure}

\section{What Computer Vision Tries to Recover}

\begin{definitionbox}{Computer Vision}
Given one image or several images, computer vision extracts properties of the three-dimensional world represented by those images.
\end{definitionbox}

A traffic scene illustrates the range of possible outputs. From the same image, a system may be asked to determine the number and type of vehicles, the location of the closest obstacle, the level of congestion, or the location and nature of the observed scene.

\begin{examplebox}{Scene understanding from a single photograph}
A visual scene may support several simultaneous questions: where the image was captured, how many people are present, what they are doing, what objects are supporting or surrounding them, and whether the situation contains an unusual or humorous interaction. These questions require more than measuring isolated pixels.
\end{examplebox}

\section{Why Visual Data Matter}

Images and video occur in personal photo collections, movies, news, sport, surveillance and security, medical imaging, scientific imaging, and many other settings. Their abundance makes visual analysis useful in both everyday systems and specialized scientific applications.

The broad motivation for computer vision includes safety, health technology, security, comfort, entertainment, and accessibility. The same field therefore spans both measurement-oriented problems and semantic interpretation.

\begin{figure}[!htbp]
\centering
\fcolorbox{LightGrey}{white}{\lectureimage{0.84\textwidth}{visual_data_domains_tight.jpg}}
\caption{Representative domains in which images and video occur: personal media, movies and news, surveillance, and scientific or medical imaging.}
\label{fig:visual-data-domains}
\end{figure}

\section{Vision and Graphics}

Computer graphics and computer vision address opposite directions of the same relationship. Graphics synthesizes images from a model of the world; vision analyzes images to infer properties of the world.

\begin{keyconcept}[Vision and graphics: opposite directions]
\centering
\begin{tikzpicture}[
  box/.style={draw=LightGrey,fill=white,rounded corners=1mm,minimum width=3.7cm,minimum height=0.82cm,align=center,font=\small\sffamily},
  arr/.style={-{Latex[length=1.8mm]},line width=0.7pt}
]
\node[box] (world) {3D world / scene model};
\node[box, right=3.9cm of world] (img) {2D image measurements};
\draw[arr,SapienzaRed] ([yshift=2.1mm]world.east) -- node[above,font=\small\sffamily]{synthesis / graphics} ([yshift=2.1mm]img.west);
\draw[arr,DeepGreen] ([yshift=-2.1mm]img.west) -- node[below,font=\small\sffamily]{analysis / vision} ([yshift=-2.1mm]world.east);
\end{tikzpicture}
\par\vspace{0.35em}
\RaggedRight The central computer-vision problem is inverse: a three-dimensional world produces a two-dimensional image, but the system must reason from the image back toward properties of that world.
\end{keyconcept}

\FloatBarrier
\begin{chaptersummary}
Computer vision converts visual measurements into information about the world. It supports both geometric measurement and semantic interpretation, and it is closely related to computer graphics through the analysis--synthesis relationship.
\end{chaptersummary}

% =============================================================================
\chapter{Why Vision Is Difficult}
% =============================================================================

Vision is difficult because an image is an incomplete and ambiguous measurement. The world is three-dimensional, images are two-dimensional, and the appearance of the same object can vary dramatically with conditions.

\section{An Ill-Posed Inverse Problem}

The real world is much more complex than what can be measured in an image. Projection from 3D to 2D removes information, so image formation cannot simply be inverted. A vision system must infer hidden properties from evidence that is incomplete and often ambiguous.

\section{Nuisance Parameters}

Several factors change image appearance without necessarily changing the identity of the object or scene.

\begin{table}[!htbp]
\centering
\caption{Major sources of variation in visual appearance.}
\label{tab:nuisance}
\rowcolors{2}{white}{PaleGrey}
\begin{tabularx}{\textwidth}{>{\sffamily\bfseries}p{0.24\textwidth}Y}
\rowcolor{TableTint}
\textbf{Factor} & \textbf{Effect on the observed image} \\
\toprule
Illumination & The same surface can appear different when lighting changes. \\
Object pose & Articulated or movable objects can change configuration. \\
Clutter & Irrelevant surrounding structure can interfere with interpretation. \\
Occlusion & Part of an object can be hidden by another object. \\
Intra-class appearance & Different members of the same category can look different. \\
Viewpoint & Appearance changes as the camera observes the object from a different direction. \\
Scale & The same object occupies different numbers of pixels at different distances or resolutions. \\
Deformation & Non-rigid objects may change shape. \\
Background clutter & Object evidence must be separated from a complicated background. \\
\bottomrule
\end{tabularx}
\end{table}

\begin{figure}[!htbp]
\centering
\fcolorbox{LightGrey}{white}{\lectureimage{0.82\textwidth}{nuisance_parameters_tight.jpg}}
\caption{The same object category can change appearance because of illumination, pose, clutter, occlusion, intra-class differences, and viewpoint.}
\label{fig:nuisance-parameters}
\end{figure}

\section{Context and Intra-Class Variation}

Intra-class variation means that objects belonging to the same category may have very different visual appearances. Context is also important: a local region that is ambiguous in isolation can become easier to interpret when its surroundings are considered.

Viewpoint, illumination, scale, deformation, occlusion, and background clutter are therefore not exceptional cases. They are normal conditions that a useful vision system must tolerate.

\section{Visual Cues}

Images are difficult to interpret, but they also reveal structure through visual cues. The task is to exploit those cues rather than treating all variation as noise.

\subsection{Depth and Ordering Cues}

Linear perspective and aerial perspective provide information about depth. Occlusion provides depth ordering: when one object hides part of another, the occluding object is interpreted as being in front.

\subsection{Shape and Lighting Cues}

Texture gradients can support shape interpretation. Shading and cast shadows also provide information about shape, position, and illumination.

\subsection{Grouping Cues}

Similarity in color, texture, or proximity can group image elements into coherent structures. Common fate provides another grouping cue: image elements moving together are likely to belong together.

\begin{figure}[!htbp]
\centering
\fcolorbox{LightGrey}{white}{\lectureimage[0.40\textheight]{0.95\textwidth}{visual_cues_grid_tight.jpg}}
\caption{Examples of visual cues: linear perspective, aerial perspective, texture gradients, and shading.}
\label{fig:visual-cues-examples}
\end{figure}


\section{Scale and Complexity}
% Keep the short chapter summary on this page rather than stranding its last line.
\enlargethispage{\baselineskip}

The computational scale of vision is large. A single image may contain thousands to millions of pixels. The space of human-recognizable object categories was described as roughly 3,000--30,000 categories, and articulated objects such as humans can involve more than 30 degrees of freedom. Visual data on the web add images, animated images, and video at very large scale.

A biological comparison reinforces the difficulty: approximately half of the cerebral cortex in primates has been associated with visual processing in the cited neurophysiological literature.

\FloatBarrier
\begin{chaptersummary}
Vision is hard because the mapping from the world to an image is many-to-one and because appearance changes with nuisance parameters. Depth, shape, lighting, and grouping cues provide information that helps recover structure despite this ambiguity.
\end{chaptersummary}

% =============================================================================
\addtocontents{toc}{\protect\setcounter{tocdepth}{0}}
\chapter{Historical Development and Changing Emphasis}
% =============================================================================

The history of computer vision shows a recurring shift between geometry, statistical modeling, data, and learning. Early optimism was followed by increasingly sophisticated representations and a recognition that visual understanding is not a short-term engineering problem.

\begin{figure}[!htbp]
\centering
\begin{tikzpicture}[font=\small\sffamily]
\draw[LightGrey,line width=1.2pt] (0,0.3)--(0,-5.7);
\foreach \y/\year/\title/\detail in {
0/1960s/{Early ambition}/{Recover structured 3D descriptions from images},
-0.9/1970s/{Structured representations}/{Representation and matching of pictorial structures},
-1.8/1980s/{Geometry and rigor}/{Mathematical formulations of vision problems},
-2.7/1990s/{Statistical recognition}/{Face recognition and statistical image analysis},
-3.6/2000s/{Broader recognition}/{Large annotated datasets and many object categories},
-4.5/2010s/{Deep learning}/{Learned representations for recognition and captioning},
-5.4/2020s/{Autonomous systems}/{Integrated perception for driving and robotics}}{
\fill[SapienzaRed] (0,\y) circle(2.4pt);
\node[anchor=east,text=SapienzaRed,font=\small\sffamily\bfseries] at (-0.3,\y) {\year};
\node[anchor=west,text width=11cm,align=left] at (0.4,\y) {\textbf{\title}\\\detail};
}
\end{tikzpicture}

\caption{Source milestones reconstructed as a readable timeline; the following sections retain the historical examples.}
\label{fig:cv-history}
\end{figure}

\section{The Early Ambition}

In 1966, Marvin Minsky assigned a summer project whose goal was essentially to connect a camera to a computer and make the machine describe what it saw. Decades later, general visual understanding remained an open research problem.

Larry Roberts' 1963 MIT PhD work, \emph{Machine Perception of Three-Dimensional Solids}, became an early landmark in interpreting synthetic scenes and recovering three-dimensional structure.

\section{1970s: Structured Representations}

Work in the 1970s made progress on selected images and structured visual representations. Fischler and Elschlager's 1973 work on the representation and matching of pictorial structures is a representative example.

\section{1980s: Geometry and Mathematical Rigor}

The 1980s saw artificial neural networks appear and then recede in influence while the field moved strongly toward geometry and more mathematically rigorous formulations.

\section{1990s: Statistical Recognition}

Face recognition became a prominent task, and statistical analysis became increasingly important in visual recognition.

\section{2000s: Broader Recognition and Data}

Recognition expanded to a broader set of categories. Large annotated datasets became increasingly important, and video processing grew as a major area of work.

\section{2010s: Deep Learning}

Deep learning returned to prominence. Representative milestones listed in the historical timeline include AlexNet (NIPS 2012), DeepFace (CVPR 2014), DeepPose (CVPR 2014), and \emph{Show, Attend and Tell} (ICML 2015).

\section{2020s: Autonomous Systems}

Autonomous vehicles represent a major modern application direction in which recognition, geometry, motion, and real-time decision making must operate together.

\FloatBarrier
\begin{chaptersummary}
Computer vision evolved from early synthetic-scene interpretation through structured representations, geometric rigor, statistical recognition, large annotated datasets, deep learning, and autonomous systems.
\end{chaptersummary}

% =============================================================================
\addtocontents{toc}{\protect\setcounter{tocdepth}{1}}
\chapter{Measurement and Semantic Understanding}
% =============================================================================

Computer vision is useful both as a measurement technology and as a way to obtain semantic information. These two goals are complementary: one emphasizes geometry and physical structure, while the other emphasizes categories and meanings that matter to people.

\section{Vision as a Measurement Device}

Measurement-oriented vision includes real-time stereo, structure from motion, multi-view stereo, and three-dimensional reconstruction from collections of images. Space exploration provides a concrete example: visual systems can support terrain reconstruction, obstacle detection, and navigation.

\begin{keyconcept}[Why learning still matters]
For some physical measurements, other sensors can be more direct or accurate than cameras. However, visual navigation and semantic tasks still require learned or data-driven interpretation, because the goals of vision include both \emph{what} is present and \emph{where} it is.
\end{keyconcept}

\section{Vision as a Source of Semantic Information}

A street scene can support object labels, scene context, and qualitative spatial information.

\begin{center}
\begin{minipage}[t]{0.48\textwidth}
  \vspace{0pt}% Align the top edge with the adjacent table, not its first baseline.
  \centering
  \fcolorbox{LightGrey}{white}{\lectureimage{\dimexpr\linewidth-2\fboxsep-2\fboxrule\relax}{semantic_object_labels_tight.jpg}}
  \captionsetup{font=footnotesize,hypcap=false}\captionof{figure}{Semantic interpretation can assign object and scene labels to visually meaningful regions.}
  \label{fig:semantic-labels}
\end{minipage}%
\hfill
\begin{minipage}[t]{0.48\textwidth}
  \vspace{0pt}
  \captionsetup{type=table,font=footnotesize,hypcap=false}
  \captionof{table}{Examples of semantic information obtainable from one scene.}
  \label{tab:semantic-levels}
  \footnotesize
  \renewcommand{\arraystretch}{1.12}
  \begin{tabularx}{\linewidth}{>{\sffamily\bfseries\RaggedRight\arraybackslash}p{0.34\linewidth}Y}
    \toprule
    \rowcolor{TableTint}
    \textbf{Level} & \textbf{Examples} \\
    \midrule
    Object categorization & sky, building, flag, face, banner, bus, wall, street lamp, cars \\
    Scene / context & outdoor, city, traffic \\
    Qualitative spatial information & horizontal, vertical, slanted, rigid moving object, non-rigid moving object \\
    \bottomrule
  \end{tabularx}
\end{minipage}
\end{center}

Semantic vision therefore includes scene type, spatial organization, and motion-related interpretation, not only isolated object labels. \textcolor{SapienzaRed}{\sffamily\bfseries Chapter summary.} Vision combines geometry and semantics: systems may need both \emph{what} is present and \emph{where} it is.

% =============================================================================
\chapter{Core Areas of Computer Vision}
% =============================================================================

A useful organization of the field separates early vision, mid-level vision, multi-view geometry, recognition, and high-level or advanced topics. The categories are not completely independent, but they provide a practical map of the techniques studied in computer vision.

\section{Early Vision: Image Formation and Processing}

Early vision begins with image formation and basic image processing. Its topics include cameras, light, color, filtering, and edges.

\subsection{The Physics of Imaging}

Camera modeling addresses what a camera does and how camera position can be inferred. Light is studied in terms of measurement, interaction with surfaces, and the way image brightness is produced. Color is studied through its underlying mechanisms and through methods for describing and measuring it.

The broader topic list includes imaging geometry, pinhole cameras, lenses, camera parameters, and calibration.

\subsection{Representing Small Image Patches}

Small image neighborhoods are represented for three main reasons:

\begin{itemize}
  \item to establish correspondence between points in different images;
  \item to detect sharp changes, commonly called edges;
  \item to represent texture through statistics of local image patterns.
\end{itemize}

Filter outputs can be interpreted as a dot-product-like comparison between a pattern and the image while the pattern is shifted across image locations. A strong response indicates that the local image resembles the pattern. Derivative-like filters are useful for detecting intensity changes.

\begin{figure}[!htbp]
\centering
\begin{tikzpicture}[
  cell/.style={draw=LightGrey,minimum width=0.42cm,minimum height=0.42cm},
  label/.style={font=\small\sffamily,align=center},
  arr/.style={-{Latex[length=2mm]},line width=0.75pt,draw=MidGrey}
]
\node[label] (im) {image};
\node[label,right=3.2cm of im] (ker) {kernel};
\node[label,right=3.2cm of ker] (out) {filter response};
\draw[arr] ([xshift=5mm]im.east) -- ([xshift=-5mm]ker.west);
\draw[arr] ([xshift=5mm]ker.east) -- ([xshift=-5mm]out.west);
\node[below=3mm of ker,font=\footnotesize\sffamily\color{MidGrey}] {shift and compare local patterns};
\end{tikzpicture}
\caption{Conceptual filtering pipeline: a kernel is shifted across the image to produce local responses.}
\label{fig:filtering-concept}
\end{figure}

Image filtering and enhancement also include linear filters, convolution, smoothing, edge detection, and image pyramids.

\begin{figure}[!htbp]
\centering
\fcolorbox{LightGrey}{white}{\lectureimage{0.93\textwidth}{early_vision_grid_tight.jpg}}
\caption{Early-vision examples: local filtering, texture classification from filter responses, and shape information derived from texture.}
\label{fig:early-vision-examples}
\end{figure}

\subsection{Texture}

Many objects are distinguished by texture. Tigers, cheetahs, grass, and trees provide intuitive examples. Texture can be represented through statistics of filter outputs: coarse bar filters respond strongly to tiger-like stripes, spot filters to cheetah-like patterns, long narrow filters to grass-like structure, and extended spot-like responses to leaves.

Objects with different textures can be segmented, and variation in texture can provide a cue to three-dimensional shape.

\section{Mid-Level Vision: Grouping, Fitting, and Tracking}

Mid-level vision seeks coherent structure by grouping or fitting lower-level measurements. Fitting can include alignment and least-squares methods.

\subsection{Segmentation}

Segmentation breaks an image or video into useful pieces. Examples include finding video sequences corresponding to one shot and finding image components that are coherent in internal appearance.

\subsection{Tracking}

Tracking keeps track of a moving object through a long sequence of views.

\Needspace{5\baselineskip}
\section{Multi-View Geometry}

\begin{figure}[!htbp]
\centering
\begin{minipage}[c]{0.45\textwidth}
\RaggedRight
Multi-view geometry studies relationships across several views. The core examples are stereo and projective structure from motion.
\end{minipage}\hfill
\begin{minipage}[c]{0.50\textwidth}
\centering
\fcolorbox{LightGrey}{white}{\lectureimage[0.34\textheight]{0.85\textwidth}{stereo_geometry_tight.jpg}}
\end{minipage}
\caption{Stereo vision uses multiple views to recover geometric information that is not directly available from a single image.}
\label{fig:stereo-geometry}
\end{figure}

\Needspace{10\baselineskip}
\section{Recognition}

A representative recognition pipeline consists of patch description and matching, clustering and visual vocabularies, bag-of-features models, and classification.

\begin{figure}[!htbp]
\centering
\begin{tikzpicture}[
  box/.style={draw=LightGrey,fill=PaleGrey,rounded corners=1mm,minimum width=3.0cm,minimum height=1.0cm,align=center,font=\small\sffamily},
  arr/.style={-{Latex[length=1.9mm]},line width=0.75pt,draw=SapienzaRed}
]
\node[box] (p) {Patch description\\and matching};
\node[box,right=7mm of p] (c) {Clustering and\\visual vocabularies};
\node[box,right=7mm of c] (b) {Bag-of-features\\models};
\node[box,right=7mm of b] (k) {Classification};
\draw[arr] (p) -- (c);
\draw[arr] (c) -- (b);
\draw[arr] (b) -- (k);
\end{tikzpicture}
\caption{Recognition pipeline from local patch description to classification.}
\label{fig:recognition-pipeline}
\end{figure}

\section{High-Level Geometric Vision}

High-level geometric vision studies relations between object geometry and image geometry. Its topics include model-based vision for finding the position and orientation of known objects, smooth surfaces and outlines, aspect graphs that describe how outlines change with viewpoint, and range data.

\section{Advanced Topics}

Representative advanced topics include segmentation, face detection, and articulated models. The broader topic list also includes image restoration, shape analysis, motion and optical flow, semantic detection, and image mosaics.

\section{Image Collections and Image-Based Rendering}

Two recurring application directions are finding images in large collections and image-based rendering. Image collections support searching and browsing. Image-based rendering addresses the difficulty of building geometric models that reproduce every real-world detail, such as surface weathering, by producing new pictures from existing images.

\FloatBarrier
\begin{chaptersummary}
Early vision studies image formation and local measurements. Mid-level vision groups and fits those measurements. Multi-view geometry reasons across viewpoints, recognition assigns categories, and high-level topics connect image evidence to object geometry and complex applications.
\end{chaptersummary}

% =============================================================================
\addtocontents{toc}{\protect\setcounter{tocdepth}{0}}
\chapter{Applications of Computer Vision}
% =============================================================================

Computer vision appears in systems ranging from document analysis to robotics and medical imaging. The examples below illustrate how the same core ideas recur in different practical settings.

\section{Detection and Recognition}

A basic detection task asks whether an image contains a person. The positive and negative examples used for this problem emphasize the need to distinguish target instances from images without the target.

\begin{figure}[!htbp]
\centering
\fcolorbox{LightGrey}{white}{\lectureimage{0.86\textwidth}{recognition_examples_tight.jpg}}
\caption{Positive and negative examples illustrate the basic structure of a visual detection problem.}
\label{fig:detection-instances}
\end{figure}

Face detection in digital cameras is a successful example: detected faces can be used to control autofocus, exposure, and flash output.

Smile detection and face recognition are further examples in which visual information is used to infer a specific facial state or identity.

\section{Optical Character Recognition}

Optical character recognition converts scanned documents into text. Related examples include handwritten digit recognition and automatic license-plate reading.

\section{Retail and Mobile Recognition}

Retail vision can detect and recognize products. A historical example is LaneHawk, in which a camera mounted at a checkout lane detects items under a shopping basket so that the cashier can verify the detected quantity.

Mobile-phone recognition extends the same idea to camera-equipped phones. Examples included Microsoft Research systems and Nokia Point \& Find.

\section{Biometrics and Password-Free Login}

Vision-based biometrics includes face, iris, and fingerprint recognition. The identification of the ``Afghan Girl'' from iris patterns is an illustrative example of iris recognition. Fingerprint scanners and face-recognition systems also support login without a conventional password.

\section{Special Effects and Interaction}

Shape capture and motion capture are used in visual effects. Examples include work associated with \emph{The Matrix} and \emph{Pirates of the Caribbean}.

Vision-based interaction includes camera-based tracking for games and interfaces. Examples include infrared tracking with the Nintendo Wii and camera tracking of crowds for interactive displays.

\section{Vehicles, Robotics, and Space Exploration}

Applications include autonomous cars, planes, missiles, robots, space exploration, and manufacturing or quality control. Smart-car systems use vision for road-scene interpretation.

A historical automotive example cited Mobileye systems in high-end BMW, GM, and Volvo models and projected broad manufacturer adoption by 2010.

Mars exploration provides a detailed robotics example. Vision systems at JPL were used for panorama stitching, three-dimensional terrain modeling, obstacle detection, and position tracking. Robotics examples also include the Mars Spirit Rover and RoboCup.

\section{Medical and Scientific Imaging}

Medical imaging includes three-dimensional imaging with modalities such as MRI and CT, as well as image-guided surgery.

\section{Image Retrieval and Three-Dimensional Viewing}

Image retrieval supports finding and browsing pictures in large collections. Three-dimensional Earth viewers are examples of large-scale three-dimensional modeling and visualization.

\Needspace{5\baselineskip}
\section{Application Overview}

\begin{figure}[!htbp]
\centering
\fcolorbox{LightGrey}{white}{\lectureimage[0.40\textheight]{0.98\textwidth}{applications_grid_tight.jpg}}
\caption{Examples of computer-vision applications in OCR, biometrics, intelligent vehicles, and medical imaging.}
\label{fig:application-examples}
\end{figure}

\begin{table}[!htbp]
\centering
\caption{Representative computer-vision application areas.}
\label{tab:applications}
\small
\renewcommand{\arraystretch}{1.08}
\begin{tabularx}{\textwidth}{Y Y}
\toprule
\rowcolor{TableTint}
\textbf{Application area} & \textbf{Representative examples} \\
\midrule
Transportation and autonomy & smart cars, autonomous vehicles, planes, robots \\
Space exploration & panorama stitching, 3D terrain modeling, obstacle detection, tracking \\
Accessibility & aid to the blind, ASL recognition \\
Manufacturing & quality control, retail product recognition \\
Security and identity & surveillance, biometrics, face / iris / fingerprint recognition \\
Documents & OCR, digit recognition, license-plate reading \\
Medical imaging & MRI, CT, 3D imaging, image-guided surgery \\
Entertainment and interaction & special effects, motion capture, camera-based interaction \\
Image collections & retrieval, browsing, image-based rendering \\
\bottomrule
\end{tabularx}
\end{table}

\FloatBarrier
\begin{chaptersummary}
Computer vision supports recognition, documents, retail, biometrics, special effects, interaction, autonomous systems, robotics, space exploration, medical imaging, and image retrieval. These applications combine different subsets of sensing, geometry, grouping, motion, and recognition.
\end{chaptersummary}

% =============================================================================
\addtocontents{toc}{\protect\setcounter{tocdepth}{1}}
\chapter{Computer Vision and Machine Learning}
% =============================================================================

Machine learning is an important tool for computer vision, but the two fields are not identical. Computer vision also includes sensing, image formation, geometric modeling, and visual representation.

\section{The Role of Machine Learning}

Machine learning studies algorithms that can change behavior based on data. A major goal is to learn complex patterns and make decisions from data. This is directly useful in computer vision, especially for recognition.

Computer vision, however, is more than learning. It also includes explicit modeling and example-based learning, and it cares about how visual data are obtained and represented.

\begin{table}[!htbp]
\centering
\caption{Contrast between the emphases described for machine learning and computer vision.}
\label{tab:cv-ml}
\begin{tabularx}{\textwidth}{>{\sffamily\bfseries}p{0.22\textwidth}Y Y}
\toprule
\rowcolor{TableTint}
 & \textbf{Machine learning emphasis} & \textbf{Computer vision emphasis} \\
\midrule
Data & Learns patterns and decisions from data & Also considers how visual data are acquired \\
Sensors & Often treats the data source as given & Includes sensor design and active vision \\
Representation & Learns or uses representations for prediction & Explicitly studies visual representation \\
Modeling & Data-driven methods are central & Modeling and learning both matter \\
\bottomrule
\end{tabularx}
\end{table}


The source's pattern-recognition example separates \textbf{training} from \textbf{testing}. Selected measurements (features) from training images are used to train a statistical method; corresponding measurements from testing images are then used to predict a category. Representative methods shown are support vector machines (SVM), Gaussian Naive Bayes (GNB), and neural networks.
\begin{figure}[!htbp]
\centering
\begin{tikzpicture}[node distance=9mm]
\node[cvnode] (train) {Training data};
\node[cvnode,right=of train] (feat) {Selected\\measurements};
\node[cvnode,right=of feat] (method) {Statistical pattern\\recognition method};
\node[cvnode,below=of train] (test) {Testing data};
\node[cvnode,right=of test] (tf) {Same selected\\measurements};
\node[cvnode,right=of tf] (pred) {Prediction};
\draw[cvarrow] (train)--(feat); \draw[cvarrow] (feat)--node[above,font=\small]{train}(method);
\draw[cvarrow] (test)--(tf); \draw[cvarrow] (tf)--(pred);
\draw[cvarrow] (method)--node[right,font=\small]{apply}(pred);
\end{tikzpicture}
\caption{Training and testing in the lecture's statistical pattern-recognition example.}
\label{fig:pattern-recognition}
\end{figure}

\section{Related Disciplines}

Computer vision is connected to image processing, pattern recognition, photogrammetry, computer graphics, artificial intelligence, machine learning, projective geometry, and control theory.

\section{Active Research Topics}

Representative research topics include object recognition, human behavior analysis, continuous computer-vision streams, biometrics and soft biometrics, large-scale three-dimensional reconstruction at city scale, medical image processing, and vision for robotics.

\section{Research Publications}

Major publication venues include \emph{IEEE Transactions on Pattern Analysis and Machine Intelligence} (TPAMI), the \emph{International Journal of Computer Vision} (IJCV), and \emph{IEEE Transactions on Image Processing}. Major conferences include ICCV, CVPR, and ECCV.

\begin{remarkbox}{Historical publication snapshot}
A historical publication snapshot lists TPAMI with an impact factor of 17.730 and IJCV with 11.541. These values are historical rather than current metrics. CVPR is described as annual, while ICCV and ECCV are described as occurring every two years.
\end{remarkbox}

The publication rankings are used to emphasize the importance of computer vision research within computer science and electrical engineering.

\section{Further Reading}

Recommended and optional books include:

\begin{itemize}
  \item David Forsyth and Jean Ponce, \emph{Computer Vision: A Modern Approach};
  \item Richard Szeliski, \emph{Computer Vision: Algorithms and Applications};
  \item Emanuele Trucco and Alessandro Verri, \emph{Introductory Techniques for 3-D Computer Vision};
  \item Rafael C. Gonzalez and Richard E. Woods, \emph{Digital Image Processing}.
\end{itemize}

\FloatBarrier
\begin{chaptersummary}
Machine learning provides important methods for visual recognition, but computer vision additionally studies cameras, sensors, image formation, geometry, representation, and visual data acquisition. The field connects to several neighboring disciplines and remains active across recognition, reconstruction, biometrics, medicine, and robotics.
\end{chaptersummary}

% =============================================================================
\clearpage
\vspace*{0.2em}
\phantomsection
\addcontentsline{toc}{chapter}{Lecture 1 Quick Review}
\markboth{Lecture 1 Quick Review}{Lecture 1 Quick Review}
\begin{center}
  {\sffamily\fontsize{22}{26}\selectfont\bfseries\color{Ink}Lecture 1 Quick Review\par}
  \vspace{0.35em}
  {\color{SapienzaRed}\rule{0.82\textwidth}{0.8pt}}
\end{center}
\vspace{0.25em}

\begin{multicols}{2}
\normalsize
\setlength{\columnsep}{0.8cm}
\setlength{\emergencystretch}{1.5em}
\reviewgroup{Foundations and visual cues}
\begin{itemize}
\item \textbf{Purpose}: extract geometric, spatial, and semantic properties of the 3D world from images; graphics follows the synthesis direction.
\item \textbf{Ambiguity}: 3D-to-2D projection is ill-posed. Illumination, pose, occlusion, viewpoint, scale, deformation, clutter, and intra-class variation affect appearance; context helps interpretation.
\item \textbf{Depth and shape}: linear/aerial perspective, occlusion, texture gradients, shading, and cast shadows provide different cues.
\item \textbf{Grouping and scale}: similarity and common fate help organization; millions of pixels and many categories make vision computationally demanding.
\end{itemize}
\reviewgroup{History and understanding}
\begin{itemize}
\item Development moved through structured representations, geometry, statistics, annotated datasets, and deep learning.
\item \textbf{Measurement} includes stereo, structure from motion, and multi-view reconstruction; \textbf{semantics} includes objects, scenes, and qualitative spatial relations.
\end{itemize}
\reviewgroup{Core areas and methods}
\begin{itemize}
\item \textbf{Early vision}: image formation, cameras, light, color, filtering, edges, and local patches. Shifted filters compare image patches with patterns.
\item \textbf{Texture}: filter-response statistics support segmentation and shape interpretation; patches also support correspondence.
\item \textbf{Mid-level and multi-view}: fitting, grouping, segmentation, tracking, stereo, and projective structure from motion.
\item \textbf{Recognition}: patch descriptions, visual vocabularies, bag-of-features, and classification. Higher geometric vision studies models, surfaces, outlines, aspect graphs, and range data.
\end{itemize}
\reviewgroup{Applications and learning}
\begin{itemize}
\item OCR, biometrics, medicine, autonomous systems, robotics, space, effects, interaction, and retrieval combine sensing, geometry, and recognition.
\item Face detection can control autofocus, exposure, and flash. Rover vision supports panoramas, terrain models, obstacles, and position tracking.
\item \textbf{Machine learning} learns patterns from data; vision additionally studies sensing, acquisition, modeling, and representation.
\item Related fields include image processing, pattern recognition, photogrammetry, graphics, AI, geometry, learning, and control.
\item Research includes recognition, behavior, biometrics, reconstruction, medicine, and robotics. Major venues include TPAMI, IJCV, IEEE TIP, CVPR, ICCV, and ECCV.
\end{itemize}
\end{multicols}


% =============================================================================
% LECTURE 2 --- IMAGE ACQUISITION
% =============================================================================
\lecturedivider{2}{Image Acquisition and Digitalization}{From scene illumination to a sampled digital image.}{Image formation $\cdot$ Digitalization $\cdot$ Pixel relations $\cdot$ Interpolation}

\chapter{Image Formation and Camera Acquisition}

The acquisition lecture follows the physical path from illumination to a digital image. A scene point receives incident light, the surface modifies that light according to its reflectance and other optical properties, the reflected light is projected onto the image plane, and the sensor converts wavelength-dependent energy into one or more measured image channels.


\begin{figure}[!htbp]
\centering
\begin{tikzpicture}[node distance=6mm]
\node[cvnode] (e) {Illumination\\$E(x,y,z,\lambda)$};
\node[cvnode,right=of e] (r) {Object / surface\\interaction $r$};
\node[cvnode,right=of r] (c) {Reflected spectral\\signal $c=Er$};
\node[cvnode,right=of c] (p) {Camera\\projection $\mathcal P$};
\node[cvnode,below=9mm of p] (cp) {Sensor plane\\$c_p(x',y',\lambda)$};
\node[cvnode,left=of cp] (v) {Spectral sensitivity\\$V(\lambda)$};
\node[cvnode,left=of v] (f) {Measured image\\$f=\int c_pV\,d\lambda$};
\node[cvnode,left=of f] (i) {Sampling and\\quantization $\to I$};
\foreach \a/\b in {e/r,r/c,c/p,p/cp,cp/v,v/f,f/i}{\draw[cvarrow] (\a)--(\b);}
\end{tikzpicture}
\caption{Acquisition chain connecting the lecture's physical, geometric, spectral, and digital stages. Symbols are defined in the following equations; $I$ is the digital image.}
\label{fig:acquisition-chain}
\end{figure}

\section{Illumination and Surface Properties}

The slides contrast typical source strengths and surface reflectivities to emphasize that observed brightness depends on both illumination and the object. Representative values shown in the lecture are collected in Table~\ref{tab:illumination-reflectivity}.

\begin{table}[htbp]
\centering
\caption{Representative illumination and reflectivity values reported in the lecture.}
\label{tab:illumination-reflectivity}
\begin{tabularx}{0.86\textwidth}{@{}Y r Y r@{}}
\toprule
\rowcolor{TableTint}
\textbf{Light source} & \textbf{$i$ (candela)} & \textbf{Object / surface} & \textbf{$r$} \\
\midrule
Solar light & 9000 & Black velvet & 0.01 \\
Cloudy sky & 1000 & White wall & 0.80 \\
Moon light & 0.01 & Silver and other metals & 0.90 \\
Internal office light & 100 & Snow & 0.93 \\
\bottomrule
\end{tabularx}
\end{table}

The lecture distinguishes five surface interactions: \textbf{absorption}, \textbf{diffusion}, \textbf{reflection}, \textbf{transparency}, and \textbf{refraction}. They describe different ways in which incoming light can be attenuated, scattered, redirected, transmitted, or bent at an interface. The wavelength $\lambda$ matters because both illumination and material response can vary spectrally.

\section{A Spectral Image-Formation Model}

For a natural image, let $E(x,y,z,\lambda)$ denote the incident light at a world point and let $r(x,y,z,\lambda)$ denote the point's reflectivity. The reflected light is modeled in the lecture as
\begin{equation}
  c(x,y,z,\lambda)=E(x,y,z,\lambda)\,r(x,y,z,\lambda).
  \label{eq:reflected-light}
\end{equation}
\mathlegend{
  \item $E(x,y,z,\lambda)$: incident spectral light at world location $(x,y,z)$;
  \item $r(x,y,z,\lambda)$: wavelength-dependent surface reflectivity;
  \item $c(x,y,z,\lambda)$: reflected spectral light leaving that point toward the imaging system;
  \item $\lambda$: wavelength.
}

\begin{keyconcept}[From scene radiance to an image]
The camera does not directly store the three-dimensional spectral function. It first projects world coordinates into two-dimensional image coordinates and then applies the wavelength sensitivity of the capture device.
\end{keyconcept}

The projection operator $\mathcal{P}$ maps a world point to camera coordinates:
\begin{equation}
  c_p(x',y',\lambda)=\mathcal{P}\!\left(c(x,y,z,\lambda)\right).
  \label{eq:projection-spectral}
\end{equation}
\mathlegend{
  \item $(x,y,z)$: world coordinates;
  \item $(x',y')$: image-plane or camera coordinates;
  \item $c_p$: projected spectral signal at the sensor plane;
  \item $\mathcal{P}$: camera projection operator.
}

Once the projected signal reaches the sensor, the device sensitivity $V(\lambda)$ determines how much of each wavelength contributes to the measured image value:
\begin{equation}
  f(x',y')=\int c_p(x',y',\lambda)V(\lambda)\,d\lambda.
  \label{eq:sensor-integration}
\end{equation}
\mathlegend{
  \item $V(\lambda)$: spectral sensitivity function of the capture device;
  \item $f(x',y')$: scalar image function measured at image location $(x',y')$;
  \item the integral: accumulation of wavelength-dependent energy weighted by sensor sensitivity.
}

\section{Perspective and Orthographic Projection}

The lecture compares two projection models. In a \textbf{perspective projection}, objects closer to the camera appear larger; this matches the human eye and ordinary cameras and is therefore natural, but mathematically more involved. In an \textbf{orthographic projection}, apparent object size is independent of distance; this is less realistic but simpler to analyze.

\begin{figure}[H]
\centering
\begin{tikzpicture}[>=Latex, font=\small\sffamily, node distance=8mm]
  \begin{scope}[xshift=-4.2cm]
    \node[font=\bfseries\color{SapienzaRed}] at (0,2.0) {Perspective};
    \draw[thick] (-2,-1.3)--(-2,1.3);
    \coordinate (C) at (-3.4,0);
    \fill[SapienzaRed] (C) circle (2pt) node[below] {camera};
    \draw[thick] (1.0,-0.6)--(1.0,0.6) node[above] {$O_1$};
    \draw[thick] (3.0,-0.6)--(3.0,0.6) node[above] {$O_2$};
    \draw[MidGrey] (C)--(1.0,0.6); \draw[MidGrey] (C)--(1.0,-0.6);
    \draw[MidGrey] (C)--(3.0,0.6); \draw[MidGrey] (C)--(3.0,-0.6);
    \node[align=center] at (0,-2.0) {same physical size;\\ nearer image is larger};
  \end{scope}
  \begin{scope}[xshift=4.2cm]
    \node[font=\bfseries\color{MutedBlue}] at (0,2.0) {Orthographic};
    \draw[thick] (-2,-1.3)--(-2,1.3);
    \draw[thick] (1.0,-0.6)--(1.0,0.6) node[above] {$O_1$};
    \draw[thick] (3.0,-0.6)--(3.0,0.6) node[above] {$O_2$};
    \foreach \y in {-0.6,0.6}{\draw[MidGrey,-{Latex[length=2mm]}] (3.2,\y)--(-2,\y);}
    \node[align=center] at (0,-2.0) {parallel rays;\\ image size is distance-independent};
  \end{scope}
\end{tikzpicture}
\caption{Conceptual contrast between perspective and orthographic projection.}
\label{fig:projection-types}
\end{figure}

\section{Color Sensors}

A color camera can be modeled as three wavelength-selective measurements at each image location. With red, green, and blue sensitivity functions $V_R$, $V_G$, and $V_B$, the lecture writes
\begin{align}
 f_R(x',y') &= \int c_p(x',y',\lambda)V_R(\lambda)\,d\lambda,\nonumber\\
 f_G(x',y') &= \int c_p(x',y',\lambda)V_G(\lambda)\,d\lambda,\label{eq:rgb-sensor}\\
 f_B(x',y') &= \int c_p(x',y',\lambda)V_B(\lambda)\,d\lambda.\nonumber
\end{align}
\mathlegend{
  \item $f_R,f_G,f_B$: the three measured image channels;
  \item $V_R,V_G,V_B$: red-, green-, and blue-tuned sensor sensitivity functions;
  \item $c_p$: projected spectral signal reaching the sensor.
}
The three channel values are then sufficient for ordinary RGB display devices to reconstruct a color appearance, even though the original spectrum contains much more information.

\FloatBarrier
\begin{chaptersummary}
Image formation is a chain: illumination $\rightarrow$ surface interaction $\rightarrow$ reflected spectral light $\rightarrow$ geometric projection $\rightarrow$ sensor sensitivity. Perspective and orthographic projection differ in whether apparent size changes with depth, while color acquisition uses multiple spectral sensitivity functions rather than one scalar response.
\end{chaptersummary}

\chapter{Digital Images: Sampling, Quantization, and Resolution}

\section{From a Continuous Image to a Digital Image}

Before digitalization, the projected image is continuous in both its spatial domain and its value range: $(x',y')\in\mathbb{R}^2$ and $f(x',y')\in\mathbb{R}$. A digital computer therefore needs two discretization steps.

\begin{definitionbox}{Sampling and quantization}
\textbf{Sampling} discretizes spatial coordinates, producing a finite grid of locations. \textbf{Quantization} discretizes the value measured at each sampled location, producing one of a finite set of representable levels.
\end{definitionbox}

The lecture expresses the two operations schematically as
\begin{align}
 (x',y') &\longrightarrow (x'_i,y'_j), \qquad i=0,\ldots,N-1,\quad j=0,\ldots,M-1, \label{eq:sampling-grid}\\
 f(x'_i,y'_j) &\longrightarrow \hat f(x'_i,y'_j)\in\mathbb{N}. \label{eq:quantization-map}
\end{align}
\mathlegend{
  \item $N,M$: numbers of sample positions along the two image axes;
  \item $(x'_i,y'_j)$: discrete sample location;
  \item $f$: continuous-valued image before quantization;
  \item $\hat f$: quantized digital value.
}


The sampling slides represent spatial sampling as multiplication by a finite \textbf{comb function}:
\begin{equation}
\operatorname{comb}(x',y')=
\sum_{i=0}^{N-1}\sum_{j=0}^{M-1}
\delta(x'-i\Delta_x,\,y'-j\Delta_y).
\label{eq:sampling-comb}
\end{equation}
\mathlegend{
\item $\delta$: two-dimensional Dirac impulse selecting a grid location;
\item $\Delta_x,\Delta_y$: spacing between samples along the two axes;
\item $i,j$: grid indices, with $N$ horizontal and $M$ vertical sample positions;
\item $x',y'$: continuous image-plane coordinates.}
The slide writes the operation schematically as
\begin{equation}
 f(x'_i,y'_j)=f(x',y')\times\operatorname{comb}(x',y'),
 \qquad f(x'_i,y'_j)=f(i,j)\ \text{for simplicity}.
\label{eq:comb-source}
\end{equation}
Here $x'_i=i\Delta_x$ and $y'_j=j\Delta_y$; $f(i,j)$ denotes the value at that sample, before quantization. The product on the right is an impulse representation of the sampled image; its impulse weights are the sample values on the left. This distinction makes the slide's shorthand precise without changing its sampling construction.

A digital image is therefore a multidimensional array. For a monochrome image, a pixel stores one intensity value; for a color image, it stores a vector of channel values. The lecture also notes other coordinate structures: $(x,y,z)$ for volumetric images such as CT scans and $(x,y,t)$ for video.

\section{Quantization and Bit Depth}

If each pixel value is represented by $b$ bits, the number of available quantization levels is
\begin{equation}
  N_c=2^b.
  \label{eq:quant-levels}
\end{equation}
\mathlegend{
  \item $N_c$: number of representable gray levels or colors;
  \item $b$: number of bits used for each quantized value.
}
For example, $b=8$ gives $2^8=256$ levels. Reducing the number of levels creates visible \textbf{false contours} because smooth intensity changes collapse into coarse steps.

\begin{figure}[!htbp]
\centering
\lectureimage[0.38\textheight]{0.95\textwidth}{quantization_levels.jpg}
\caption{Source examples of progressively coarser quantization; false contouring becomes evident at very low bit depth.}
\label{fig:quantization-levels}
\end{figure}

After sampling and quantization, the lecture describes an image as a discrete function with a finite discrete domain and range:
\begin{equation}
 \Omega\subset\mathbb N^2,\qquad I:\Omega\longrightarrow\{1,\ldots,P\},\qquad P\in\mathbb N.
 \label{eq:discrete-image}
\end{equation}
\mathlegend{\item $\Omega$: finite set of pixel coordinates selected by sampling;
\item $I$: image assigning a quantized value to each pixel;
\item $P$: number of representable levels; the source's 8-bit example uses $P=2^8=256$ and $\log_2 P=8$.}
The slide numbers the levels from $1$ to $P$; this expresses the number of levels, while binary image values elsewhere use $0$ and $1$.

\textbf{Digital image processing} means processing digital, multidimensional images using a digital computer.

\section{Digital Image Types}

\begin{table}[htbp]
\centering
\caption{Digital image representations presented in the lecture.}
\label{tab:image-types}
\begin{tabularx}{\textwidth}{@{}l Y Y@{}}
\toprule
\rowcolor{TableTint}
\textbf{Type} & \textbf{Pixel content} & \textbf{Interpretation} \\
\midrule
Intensity / monochrome & one scalar & light intensity, commonly displayed as gray level \\
RGB color & a 3-component vector & red, green, and blue channel values \\
Binary & one bit & $0$ for black and $1$ for white in the lecture convention \\
Indexed & one integer index & lookup into a color table containing RGB triples \\
\bottomrule
\end{tabularx}
\end{table}

\section{Pixels, Memory, and Spatial Resolution}

A \textbf{pixel} is a single value on the sampled grid. The lecture stresses that a pixel should be understood as the smallest treatable image area, not as an infinitesimal geometric point.

For an $M\times N$ image with $L$ intensity levels, binary storage uses
\begin{equation}
  k=\log_2 L, \qquad B=M N k,
  \label{eq:image-memory}
\end{equation}
\mathlegend{
  \item $M,N$: image dimensions in pixels;
  \item $L$: number of intensity levels;
  \item $k$: bits per pixel value;
  \item $B$: total number of bits for the image.
}
\begin{remarkbox}{Notation in the source}
One slide writes $k=\log(L)$, while the later 8-bit example explicitly uses $\log_2(P)=8$ for $P=256$. The binary-storage interpretation therefore uses base-2 logarithms.
\end{remarkbox}

\textbf{Spatial resolution} is the smallest distinguishable detail in an image. The lecture lists pixels per inch (ppi), dots per inch (dpi), and raw $M\times N$ dimensions as common descriptors, while warning that pixel dimensions alone do not fully express the relation between digital sampling and the detail present in the original scene.

\begin{figure}[!htbp]
\centering
\lectureimage[0.34\textheight]{0.92\textwidth}{spatial_resolution_effect.jpg}
\caption{Effect of reducing spatial sampling while preserving the same field of view: fine detail is progressively lost.}
\label{fig:spatial-resolution}
\end{figure}

Under-sampling is irreversible: interpolation can produce a denser grid, but it cannot recover scene details that were never sampled. The lecture gives application-dependent resolution examples, such as monitor/web display around 72 dpi and printing at substantially higher dpi, and presents interpolation as an approximation when increasing image size.

\begin{keyconcept}[Resolution and perceived quality]
For images displayed at the same physical size, insufficient spatial sampling produces visible pixels and lost detail. If the pixel dimensions are kept fixed instead, changing the number of samples also changes the displayed image size. The suitable combination of image size and pixel depth is application- and content-dependent.
\end{keyconcept}

\FloatBarrier
\begin{chaptersummary}
Digitalization separates two independent limitations: spatial sampling controls where the image is measured, while quantization controls how finely each measured value is represented. Bit depth determines the number of levels; spatial resolution determines recoverable detail; interpolation only estimates values between existing samples.
\end{chaptersummary}

\chapter{Pixel Relationships and Image Resampling}

\section{Neighborhoods}

For a pixel $p=(x,y)$, the lecture defines the four-neighborhood, diagonal neighborhood, and eight-neighborhood:
\begin{align}
N_4(p) &= \{(x-1,y),(x+1,y),(x,y-1),(x,y+1)\},\label{eq:n4}\\
N_D(p) &= \{(x-1,y-1),(x+1,y-1),(x-1,y+1),(x+1,y+1)\},\label{eq:nd}\\
N_8(p) &= N_4(p)\cup N_D(p).\label{eq:n8}
\end{align}
\mathlegend{
  \item $p=(x,y)$: reference pixel;
  \item $N_4(p)$: horizontal and vertical neighbors;
  \item $N_D(p)$: four diagonal neighbors;
  \item $N_8(p)$: all eight surrounding grid locations.
}

\begin{figure}[!htbp]
\centering
\begin{tikzpicture}[scale=0.8,font=\small\sffamily]
  \foreach \i in {-1,0,1}{\foreach \j in {-1,0,1}{\draw[LightGrey,fill=PaleGrey] (\i,\j) rectangle ++(1,1);}}
  \fill[SapienzaRed!20] (0,0) rectangle ++(1,1);
  \node[font=\bfseries] at (0.5,0.5) {$p$};
  \foreach \x/\y in {-1/0,1/0,0/-1,0/1}{\fill[MutedBlue!20] (\x,\y) rectangle ++(1,1);}
  \foreach \x/\y in {-1/-1,-1/1,1/-1,1/1}{\fill[DeepGreen!18] (\x,\y) rectangle ++(1,1);}
  \node[anchor=west] at (2.2,1.0) {blue: $N_4(p)$};
  \node[anchor=west] at (2.2,0.4) {green: $N_D(p)$};
  \node[anchor=west] at (2.2,-0.2) {$N_8(p)=N_4(p)\cup N_D(p)$};
\end{tikzpicture}
\caption{Pixel-neighborhood relations used to define connectivity, adjacency, and paths.}
\label{fig:pixel-neighborhoods}
\end{figure}

\section{Connectivity, Adjacency, and Paths}

Two pixels from the same class are connected when the selected neighborhood relation links them. The lecture distinguishes \textbf{4-connectivity}, \textbf{8-connectivity}, and \textbf{mixed (m-) connectivity}. In the stated m-connectivity rule, $p$ and $q$ are connected if either $q\in N_4(p)$, or $q\in N_D(p)$ and $N_4(p)\cap N_4(q)=\varnothing$. The mixed rule is introduced to avoid the ambiguity that diagonal links can create under pure 8-connectivity.

\begin{figure}[!htbp]
\centering
\begin{tikzpicture}[font=\small\sffamily,x=0.85cm,y=0.85cm]
\foreach \shift/\heading in {0/{8-path: ambiguous},5/{m-path: source illustration}}{
\begin{scope}[xshift=\shift cm]
\foreach \x/\y in {0/2,1/2,0/1,1/0}{\draw[MidGrey,fill=PaleRed] (\x,\y) rectangle ++(1,1);}
\node[above] at (1,3.4) {\heading};
\node[right] at (2,2.5) {$p$}; \node[below] at (1.5,0) {$q$};
\draw[cvarrow] (1.5,2.5)--(0.5,2.5)--(0.5,1.5)--(1.5,0.5);
\end{scope}}
\draw[SapienzaRed,dashed,line width=0.9pt] (1.5,2.5)--(0.5,1.5);
\end{tikzpicture}
\caption{The source example removes the alternative diagonal route in its m-path illustration, retaining a single route from $p$ to $q$.}
\label{fig:connectivity-paths}
\end{figure}
\begin{remarkbox}{Source rule and illustration}
The displayed rule uses $N_4(p)\cap N_4(q)=\varnothing$ without restricting the intersection to the selected pixel class. On the full interior grid, diagonal pixels share two 4-neighbors, so the literal rule excludes all such diagonal links. The source's path illustration nevertheless retains the final diagonal step. Both are preserved here: the diagram reproduces the slide's intended ambiguity removal, rather than silently replacing its written definition with another convention.
\end{remarkbox}


A pixel $p$ is \textbf{adjacent} to $q$ when they are connected. Two subsets $S_1$ and $S_2$ are adjacent if some pixel in one subset is adjacent to a pixel in the other. A path from $p=(x,y)$ to $q=(s,t)$ is a sequence of distinct adjacent pixels $(x_0,y_0),\ldots,(x_n,y_n)$ with the required endpoints; the adjacency convention determines whether the result is a 4-path, 8-path, or m-path.

\section{Distance Measures}

A distance $D$ must satisfy non-negativity, identity of indiscernibles, symmetry, and the triangle inequality. Three distances appear in the lecture:
\begin{align}
D_E(p,q) &= \sqrt{(x-s)^2+(y-t)^2},\label{eq:euclidean-distance}\\
D_4(p,q) &= |x-s|+|y-t|,\label{eq:cityblock-distance}\\
D_8(p,q) &= \max\{|x-s|,|y-t|\}.\label{eq:chessboard-distance}
\end{align}
\mathlegend{
  \item $p=(x,y)$ and $q=(s,t)$: two pixel coordinates;
  \item $D_E$: Euclidean distance;
  \item $D_4$: city-block distance, whose unit-distance pixels are exactly the 4-neighbors;
  \item $D_8$: chessboard distance, whose unit-distance pixels are exactly the 8-neighbors.
}

\section{Interpolation}

Interpolation estimates the value of an image at non-integer coordinates. The lecture separates resizing into two steps:
\begin{enumerate}
\item \textbf{Create the new sampling grid}, mapping its positions into the original image coordinates. The slide's $1.5\times$ zoom changes a $500\times500$ grid to $750\times750$.
\item \textbf{Assign a value to each new pixel} by interpolation of the existing samples.
\end{enumerate}
The lecture compares four approaches.

\begin{table}[htbp]
\centering
\caption{Interpolation methods discussed in the acquisition lecture.}
\label{tab:interpolation-methods}
\begin{tabularx}{\textwidth}{@{}l Y Y@{}}
\toprule
\rowcolor{TableTint}
\textbf{Method} & \textbf{Principle} & \textbf{Lecture emphasis} \\
\midrule
Nearest neighbor & choose the closest sampled pixel & fastest; blocky appearance; be careful with integer rounding \\
Triangle interpolation & weighted combination of three triangle vertices & useful for less-structured samples; weights follow opposite sub-areas and should be normalized \\
Bilinear & weighted combination of the four corners of a grid cell & smoother than nearest neighbor; four lookups; often a good speed/quality tradeoff \\
Bicubic & cubic interpolation using a $4\times4$ neighborhood & 16 neighbors; smoother than bilinear; more computation \\
\bottomrule
\end{tabularx}
\end{table}

Nearest-neighbor interpolation is written in the slides as
\begin{equation}
  f(x,y,z)=\operatorname{Im}(\operatorname{round}(x),\operatorname{round}(y),z),
  \label{eq:nearest-neighbor}
\end{equation}
Here $f(x,y,z)$ is the interpolated image value, $\operatorname{Im}$ is the original sampled image, $x,y$ are source spatial coordinates, $\operatorname{round}$ selects the nearest integer, and $z$ remains an integer channel index in the implementation example.

For triangle interpolation, the slide uses the schematic weighted sum
\begin{equation}
  Q=V_1A_1+V_2A_2+V_3A_3,
  \label{eq:triangle-interp}
\end{equation}
Here $Q$ is the value assigned at the interior point $q$, $V_i$ are vertex values, and $A_i$ are opposite sub-triangle areas, for $i=1,2,3$. This is the \textbf{slide's schematic expression}. The slide explicitly requires normalization by total area in the general case. Its \textbf{normalized interpretation} is
\begin{equation}
 Q=\frac{V_1A_1+V_2A_2+V_3A_3}{A_1+A_2+A_3}.
 \label{eq:triangle-normalized}
\end{equation}
The denominator is the whole triangle's area; the weights $A_i/(A_1+A_2+A_3)$ therefore sum to one.

\begin{figure}[!htbp]
\centering
\begin{tikzpicture}[font=\small,scale=1.15]
\coordinate (v1) at (0,3); \coordinate (v2) at (3,0); \coordinate (v3) at (-1,0.6);
\coordinate (q) at (0.5,1.1);
\fill[PaleGreen] (v2)--(v3)--(q)--cycle;
\fill[PaleRed] (v1)--(v3)--(q)--cycle;
\fill[PaleGrey] (v1)--(v2)--(q)--cycle;
\draw[thick,MidGrey] (v1)--(v2)--(v3)--cycle;
\foreach \v in {v1,v2,v3}{\draw[MidGrey] (q)--(\v); \fill[SapienzaRed] (\v) circle(2pt);}
\node[above] at (v1) {$V_1$}; \node[right] at (v2) {$V_2$}; \node[left] at (v3) {$V_3$};
\fill[SapienzaRed] (q) circle(2pt) node[above left] {$q$};
\node at (0.8,0.55) {$A_1$}; \node at (-0.2,1.55) {$A_2$}; \node at (1.2,1.65) {$A_3$};
\end{tikzpicture}
\caption{Triangle construction: $A_i$ is the sub-triangle opposite vertex $V_i$. Normalize by $A_1+A_2+A_3$ to obtain area weights.}
\label{fig:triangle-geometry}
\end{figure}



For bilinear interpolation inside a unit grid cell, the four opposite rectangle areas are
\begin{equation}
A_1=d_2d_4,\quad A_2=d_1d_4,\quad A_3=d_2d_3,\quad A_4=d_1d_3,
\end{equation}
and the interpolated value is
\begin{equation}
  Q=V_1A_1+V_2A_2+V_3A_3+V_4A_4.
  \label{eq:bilinear-area}
\end{equation}
\mathlegend{
  \item $Q$: interpolated value at the position $q$; $V_1,\ldots,V_4$: values at the four surrounding grid points;
  \item $d_1,d_2,d_3,d_4$: horizontal and vertical fractional distances shown in the lecture construction;
  \item $A_i$: area opposite vertex $V_i$; within a unit cell the weights sum to one.
}

\begin{figure}[!htbp]
\centering
\begin{tikzpicture}[font=\small,x=1.35cm,y=1.15cm]
\fill[PaleGreen] (0,1.1) rectangle (2.1,3);
\fill[PaleRed] (0,0) rectangle (2.1,1.1);
\fill[LightGrey!55] (2.1,1.1) rectangle (3,3);
\fill[PaleAmber] (2.1,0) rectangle (3,1.1);
\draw[thick,MidGrey] (0,0) rectangle (3,3);
\draw[MidGrey] (2.1,0)--(2.1,3); \draw[MidGrey] (0,1.1)--(3,1.1);
\foreach \x/\y/\v in {0/3/1,3/3/2,0/0/3,3/0/4}{\fill[SapienzaRed] (\x,\y) circle(2pt);}
\node[above left] at (0,3) {$V_1$}; \node[above right] at (3,3) {$V_2$};
\node[below left] at (0,0) {$V_3$}; \node[below right] at (3,0) {$V_4$};
\fill[SapienzaRed] (2.1,1.1) circle(2pt) node[above left] {$q$};
\node at (2.55,0.55) {$A_1$}; \node at (1.05,0.55) {$A_2$};
\node at (2.55,2.05) {$A_3$}; \node at (1.05,2.05) {$A_4$};
\draw[<->,MidGrey] (0,3.45)--node[above]{$d_1$}(2.1,3.45);
\draw[<->,MidGrey] (2.1,3.45)--node[above]{$d_2$}(3,3.45);
\draw[<->,MidGrey] (3.55,3)--node[right]{$d_3$}(3.55,1.1);
\draw[<->,MidGrey] (3.55,1.1)--node[right]{$d_4$}(3.55,0);
\end{tikzpicture}
\caption{Bilinear geometry redrawn from the source: $d_1+d_2=d_3+d_4=1$ in the unit cell. Each corner uses the opposite rectangle's area; $Q$ is the value at $q$. The drawing is enlarged for readability.}
\label{fig:bilinear-geometry}
\end{figure}


Bicubic interpolation is presented as a cubic interpolation of four cubic interpolations. It uses the 16 nearest samples, fits third-order polynomials such as $f(x)=a+bx+cx^2+dx^3$, interpolates along one axis, and then fits/interpolates along the other. Here $x$ is the interpolation coordinate, $f(x)$ is its estimated value, and $a,b,c,d$ are the fitted polynomial coefficients.

\begin{figure}[!htbp]
\centering
\lectureimage[0.34\textheight]{0.90\textwidth}{bicubic_vs_bilinear.jpg}
\caption{Lecture comparison of bilinear and bicubic resampling on the same image region.}
\label{fig:bicubic-vs-bilinear}
\end{figure}

The resize algorithm is consequently simple at a high level: for every pixel in the new image, map its position back into the original-image coordinates, interpolate the source value, and write that value into the output image.

\FloatBarrier
\begin{chaptersummary}
Neighborhood relations define connectivity, adjacency, paths, and grid-aware distance measures. Resampling then estimates values between grid points: nearest neighbor is cheapest, bilinear provides a common compromise, and bicubic trades additional computation for smoother results. Increasing sample density by interpolation does not reverse information lost during down-sampling.
\end{chaptersummary}

\clearpage
\phantomsection
\addcontentsline{toc}{chapter}{Lecture 2 Quick Review}
\markboth{Lecture 2 Quick Review}{Lecture 2 Quick Review}
\begin{center}
{\sffamily\fontsize{22}{26}\selectfont\bfseries\color{Ink}Lecture 2 Quick Review\par}
\vspace{0.35em}{\color{SapienzaRed}\rule{0.82\textwidth}{0.8pt}}
\end{center}

\normalsize
\reviewgroup{Physical acquisition}
\begin{enumerate}[label=\textcolor{SapienzaRed}{\bfseries\arabic*.},leftmargin=1.8em,itemsep=0.40em]
\item Reflected light is modeled by incident illumination multiplied by surface reflectivity.
\item Projection maps the 3D scene into 2D image coordinates.
\item Sensor sensitivity integrates the projected spectrum over wavelength.
\item RGB acquisition uses three different spectral sensitivity functions.
\item Perspective projection changes apparent size with depth; orthographic projection does not.
\end{enumerate}
\reviewgroup{Digitalization and storage}
\begin{enumerate}[resume,label=\textcolor{SapienzaRed}{\bfseries\arabic*.},leftmargin=1.8em,itemsep=0.40em]
\item Sampling discretizes position; quantization discretizes value.
\item $b$ bits provide $2^b$ representable levels.
\item Low quantization depth can create false contours.
\item A pixel is the smallest treatable sampled image area.
\item For $M\times N$ pixels and $k$ bits per value, storage is $MNk$ bits.
\end{enumerate}
\reviewgroup{Pixel relations and interpolation}
\begin{enumerate}[resume,label=\textcolor{SapienzaRed}{\bfseries\arabic*.},leftmargin=1.8em,itemsep=0.40em]
\item $N_4$, $N_D$, and $N_8$ define local grid neighborhoods.
\item m-connectivity avoids some diagonal ambiguity of 8-connectivity.
\item $D_4$ is city-block distance; $D_8$ is chessboard distance.
\item Nearest neighbor is fast but blocky; bilinear uses four samples; bicubic uses sixteen.
\item Interpolation creates estimates, not lost high-frequency detail.
\end{enumerate}


% =============================================================================
% LECTURE 3 --- COLOR
% =============================================================================
\lecturedivider{3}{Color, Perception, and Texture}{From spectral measurements to color and appearance.}{Human vision $\cdot$ Color spaces $\cdot$ White balance $\cdot$ Texture}

\chapter{Color: Light and Human Vision}

\section{The Nature of Color}

The lecture emphasizes that color is not simply a physical property attached to an object. It is a \textbf{psychological property of visual experience} produced by the interaction between the spectrum of physical light and the human visual system. This distinction motivates both colorimetry and color constancy: a camera measures spectral effects, whereas a human observer experiences a perceptual color.

Visible light occupies only a limited interval of the electromagnetic spectrum, approximately $400$--$700\,\mathrm{nm}$ in the lecture. A light source can be described physically by its spectral power distribution: the amount of energy emitted at each wavelength. The lecture contrasts narrow-band sources such as a ruby laser with broader spectra such as tungsten illumination and daylight.

\begin{figure}[!htbp]
\centering
\lectureimage[0.36\textheight]{0.94\textwidth}{source_spectra.jpg}
\caption{Examples from the lecture of qualitatively different source spectra, including narrow-band and broadband illumination.}
\label{fig:source-spectra}
\end{figure}

A surface is similarly characterized by its reflectance spectrum: the fraction of incident light reflected at each wavelength. Therefore the observed color signal depends jointly on source spectrum and surface reflectance.

\begin{figure}[!htbp]
\centering
\lectureimage[0.32\textheight]{0.96\textwidth}{surface_reflectance_spectra.jpg}
\caption{Representative surface reflectance spectra shown in the lecture for red, yellow, blue, and purple surfaces.}
\label{fig:surface-reflectance}
\end{figure}

\section{The Eye as a Sensing System}

The lecture draws a camera analogy for the eye: the iris regulates the pupil aperture, the lens changes shape to focus at different distances, and the retina supplies the photoreceptors. Rods and cones are not uniformly distributed. Rods are emphasized for intensity sensitivity, while cones support color vision. The \textbf{fovea}, a small central region covering roughly $1$--$2^{\circ}$ of visual angle in the slides, has the highest cone density and no rods; peripheral acuity is lower because many rods can feed the same neuron.

Human color vision uses three cone classes, conventionally called S, M, and L for their relative wavelength sensitivities. The lecture gives an approximate L:M:S ratio of $10:5:1$ and notes that there are almost no S cones in the center of the fovea.

\begin{figure}[!htbp]
\centering
\lectureimage[0.34\textheight]{0.93\textwidth}{cone_sensitivity.jpg}
\caption{Lecture illustration of the three cone-sensitivity functions and their non-uniform representation.}
\label{fig:cone-sensitivity}
\end{figure}

Rods operate at lower illumination than cones. Below the color-perception threshold, rod-dominated vision can still detect brightness, while cone-based color and fine foveal detail become ineffective. This connects the sensitivity slide's question, ``Why can't we read in the dark?'', to the absence of rods in the fovea and the lower acuity of peripheral rod vision. The source chart distinguishes the threshold in bright sunlight from the much lower threshold after dark adaptation. Light/dark adaptation changes sensitivity as illumination changes; it is a gradual effect, unlike instantaneous contrast.

\begin{figure}[!htbp]
\centering
\lectureimage[0.30\textheight]{0.82\textwidth}{rod_cone_regimes.png}
\caption{Source sensitivity regimes and light/dark adaptation thresholds. The chart is preserved rather than assigning invented numerical sensitivity curves.}
\label{fig:rod-cone-regimes}
\end{figure}

Brewster's-color demonstrations are used as evidence that the visual system can interpolate color from spatially offset samples. This is consistent with the fact that the three cone populations do not sample every retinal location identically.

\section{Spectral Filtering and Metamers}

A cone response can be understood as a spectral filter: multiply the incoming spectrum by the receptor sensitivity curve and integrate over wavelength. Each cone class therefore compresses an entire spectrum into a single response. Three cone outputs cannot preserve the full spectral signal, so distinct spectra can produce identical responses.

\begin{definitionbox}{Metamerism}
\textbf{Metamers} are physically different spectral power distributions that produce the same perceived color under the relevant viewing conditions. Metamerism is a direct consequence of compressing a high-dimensional spectrum into a small number of receptor responses.
\end{definitionbox}

\begin{figure}[!htbp]
\centering
\lectureimage[0.28\textheight]{0.86\textwidth}{metamer_spectra.jpg}
\caption{Real-world spectral curves used in the lecture to illustrate that different spectra may be perceptually matched.}
\label{fig:metamer-spectra}
\end{figure}

\FloatBarrier
\begin{chaptersummary}
Physical spectra are high-dimensional, but human color perception is driven mainly by three cone classes. Consequently, color depends on both illumination and surface reflectance, and different spectra can become perceptually indistinguishable as metamers.
\end{chaptersummary}

\chapter{Color Matching and Color Spaces}

\section{Trichromacy and Color Matching}

Color-matching experiments standardize which spectra produce the same sensation under controlled conditions. The lecture's central empirical result is \textbf{trichromacy}: for most observers, three primaries are sufficient to match a wide range of test lights, with color blindness noted as an exception. The same test light and same primaries tend to produce similar matching weights across observers.

In the experiment, the observer compares a test light with a mixture of three primary lights and adjusts their weights until the two fields appear to match.
\begin{figure}[!htbp]
\centering
\begin{tikzpicture}[node distance=8mm]
\node[cvnode] (test) {Test light};
\node[cvnode,right=23mm of test,text width=38mm] (obs) {Observer compares\\the two light fields};
\node[cvnode,right=23mm of obs] (mix) {Mixture of three\\primary lights};
\draw[cvarrow] (test)--(obs); \draw[cvarrow] (mix)--(obs);
\draw[cvarrow] (obs.south)--++(0,-0.6)-|node[pos=0.3,below,font=\small\sffamily]{adjust primary weights until a perceptual match}(mix.south);
\end{tikzpicture}
\caption{Conceptual color-matching experiment underlying trichromacy and Grassmann's laws.}
\label{fig:color-matching-experiment}
\end{figure}

\section{Grassmann's Laws}

The lecture states that additive color matching is linear. If
\begin{align*}
A &= u_1P_1+u_2P_2+u_3P_3,\\
B &= v_1P_1+v_2P_2+v_3P_3,
\end{align*}
then mixtures and scalings obey
\begin{align}
A+B &= (u_1+v_1)P_1+(u_2+v_2)P_2+(u_3+v_3)P_3,\label{eq:grassmann-add}\\
kA &= (ku_1)P_1+(ku_2)P_2+(ku_3)P_3.\label{eq:grassmann-scale}
\end{align}
\mathlegend{
  \item $P_1,P_2,P_3$: chosen primary lights;
  \item $u_i,v_i$: primary weights that match test lights $A$ and $B$;
  \item $k$: scalar intensity multiplier.
}
If two test lights are matched by the same three weights, the color-matching experiment treats them as a perceptual match.

\section{Color-Matching Functions as a Linear Transform}

Choose three primaries and measure the weights $c_1(\lambda)$, $c_2(\lambda)$, and $c_3(\lambda)$ required to match monochromatic lights at sampled wavelengths. Store these color-matching functions as the rows of a matrix
\begin{equation}
C=\begin{bmatrix}
 c_1(\lambda_1)&\cdots&c_1(\lambda_N)\\
 c_2(\lambda_1)&\cdots&c_2(\lambda_N)\\
 c_3(\lambda_1)&\cdots&c_3(\lambda_N)
\end{bmatrix}.
\end{equation}
Represent a new sampled spectral signal by
\[
\mathbf{t}=\begin{bmatrix}t(\lambda_1)&\cdots&t(\lambda_N)\end{bmatrix}^{\mathsf T}.
\]
The primary coordinates are then
\begin{equation}
  \mathbf{e}=C\mathbf{t}.
  \label{eq:color-matching-matrix}
\end{equation}
\mathlegend{
  \item $C$: $3\times N$ matrix of sampled color-matching functions;
  \item $\mathbf{t}$: sampled spectral power distribution of the test light;
  \item $\mathbf{e}=(e_1,e_2,e_3)^{\mathsf T}$: three primary weights that describe its color in the chosen linear color space; $\lambda_1,\ldots,\lambda_N$ are the $N$ sampled wavelengths, $t(\lambda_i)$ is spectral power, and $c_j(\lambda_i)$ is the matching weight for primary $j$.
}

\section{RGB and CIE XYZ}

A \textbf{linear color space} is defined by a choice of three primaries. Mixing two lights traces a straight line between their coordinates; mixtures of three primaries lie within the triangle they define. In RGB, the primaries correspond to red, green, and blue lights. The source figure specifies $p_1=645.2\,\mathrm{nm}$ (red), $p_2=525.3\,\mathrm{nm}$ (green), and $p_3=444.4\,\mathrm{nm}$ (blue); these are the selected monochromatic primaries in that lecture example. The lecture notes that some wavelengths require subtractive matching, which appears as negative values in RGB color-matching functions.

The CIE XYZ system, standardized in 1931 by the International Commission on Illumination, uses imaginary primaries chosen so that its matching functions are nonnegative. A two-dimensional chromaticity representation is obtained by normalizing:
\begin{equation}
  x=\frac{X}{X+Y+Z}, \qquad y=\frac{Y}{X+Y+Z}.
  \label{eq:cie-xy}
\end{equation}
\mathlegend{
  \item $X,Y,Z$: CIE tristimulus coordinates;
  \item $x,y$: normalized chromaticity coordinates that remove the common scale $X+Y+Z$ (defined when this sum is nonzero).
}

\begin{figure}[!htbp]
\centering
\lectureimage[0.38\textheight]{0.95\textwidth}{cie_xyz_chromaticity.jpg}
\caption{CIE XYZ material from the lecture: chromaticity visualization and the associated nonnegative matching functions.}
\label{fig:cie-xyz}
\end{figure}

\section{Toward Perceptual Uniformity}

The lecture stresses that Euclidean differences in CIE $x,y$ coordinates are not perceptually uniform. McAdam ellipses illustrate that just-noticeable color differences occupy differently sized and oriented regions of the chromaticity diagram. CIE $u'v'$ is introduced as a projective transform intended to make these regions more uniform, followed by CIE $L^*a^*b^*$ as a later perceptually motivated space.

\begin{figure}[!htbp]
\centering
\lectureimage[0.35\textheight]{0.90\textwidth}{mcdam_ellipses.jpg}
\caption{McAdam-ellipse illustration used to motivate more perceptually uniform color coordinates.}
\label{fig:mcdam}
\end{figure}

\FloatBarrier
\section{HSV}

HSV reorganizes color into perceptually interpretable dimensions: \textbf{Hue} answers ``what color?'', \textbf{Saturation} describes how much chromatic purity is present, and \textbf{Value} describes brightness. The lecture visualizes the model as a cylinder or cone and interprets it geometrically as a reorganization of the RGB cube. In the source's geometric construction, the RGB cube is tilted onto its black--white diagonal: that diagonal becomes the brightness axis, while chromatic colors lie around it. Hue is the angular direction around this axis, saturation moves outward from the neutral axis, and value moves along the brightness direction. The source then schematically places hues in a common plane, expands the cube, and smooths it toward a cylinder. This representation makes operations such as shifting hue or increasing saturation conceptually simple.

\begin{figure}[H]
\centering
\lectureimage[0.35\textheight]{0.95\textwidth}{hsv_model.jpg}
\caption{Geometric HSV representation shown in the lecture, emphasizing hue, saturation, and value as different axes of description.}
\label{fig:hsv-model}
\end{figure}

\FloatBarrier
\begin{chaptersummary}
Trichromacy makes three-coordinate color spaces possible. Grassmann's laws justify linear color matching, while $\mathbf e=C\mathbf t$ maps a sampled spectrum into primary coordinates. RGB is device-oriented, CIE XYZ provides a standardized linear basis, CIE $u'v'$ and $L^*a^*b^*$ address perceptual uniformity, and HSV reorganizes RGB-like color into perceptual dimensions.
\end{chaptersummary}

\chapter{Color Constancy and White Balance}

\section{Constancy and Context}

\textbf{Color/lightness constancy} is the ability of the human visual system to perceive relatively stable surface reflectance despite changes in illumination. The lecture separates several classes of perceptual effects:
\begin{itemize}
  \item \textbf{instantaneous}: simultaneous contrast and Mach bands;
  \item \textbf{gradual}: light/dark adaptation, chromatic adaptation, and afterimages;
  \item \textbf{high-level}: effects related to color naming and scene interpretation.
\end{itemize}

\textbf{Simultaneous contrast} means that a target's apparent lightness or color changes with its surrounding field. \textbf{Mach bands} emphasize apparent bright and dark bands near transitions between neighboring intensity regions: perceived edge contrast can exceed the physical change.

\begin{figure}[!htbp]
\centering
\lectureimage[0.24\textheight]{0.86\textwidth}{simultaneous_contrast_source.png}
\caption{The source's simultaneous-contrast/Mach-band demonstration. The colored targets are viewed within alternating surrounding stripes; the perceptual artwork is preserved.}
\label{fig:simultaneous-contrast}
\end{figure}


The checker-shadow illusion shows that equal measured image intensities can be perceived differently. A boundary between checker squares suggests a \textbf{reflectance change} (different surface material), whereas the cast-shadow boundary suggests an \textbf{illumination change} across a surface. Interpreting the shadow allows the observer to judge surface lightness rather than treating every intensity change as a change of reflectance.

\begin{figure}[!htbp]
\centering
\lectureimage[0.27\textheight]{0.72\textwidth}{checker_shadow.jpg}
\caption{Checker-shadow example from the lecture, used to discuss lightness constancy and the distinction between reflectance and illumination changes.}
\label{fig:checker-shadow}
\end{figure}

\section{Chromatic Adaptation}

The visual system changes sensitivity according to the prevailing luminance and chromatic conditions. Aperture changes regulate total light entering the eye, while receptor sensitivity also adapts. The lecture's example is that under increased red illumination, red-sensitive responses are gradually reduced until the scene appears more neutral. Adaptation works better in brighter scenes, which helps explain why candlelit scenes may still look yellow.

\begin{warningbox}[Human perception versus a computer]
Humans use context across the entire scene and perform strong perceptual normalization. A computer ordinarily starts from measured pixel values and does not automatically separate an object's intrinsic color from illumination. Algorithms such as white balance attempt to approximate part of this adaptation.
\end{warningbox}

\section{White Balance}

When a photograph is viewed on a screen or print, the observer adapts to the room's illuminant rather than the illuminant present in the photographed scene. Incorrect camera white balance therefore appears as an unnatural color cast.

\begin{figure}[!htbp]
\centering
\lectureimage[0.30\textheight]{0.94\textwidth}{white_balance_example.jpg}
\caption{Incorrect and corrected white balance as illustrated in the lecture.}
\label{fig:white-balance-example}
\end{figure}

The lecture contrasts film and digital practice. Film cameras use different films or filters for different illuminants, whereas digital cameras can use automatic white balance, predefined illuminant settings, or a custom neutral reference. Multiple illuminants make the task ambiguous because different reference regions can imply different corrections.

\begin{figure}[!htbp]
\centering
\lectureimage[0.30\textheight]{0.90\textwidth}{white_balance_mixed_illuminants.jpg}
\caption{Mixed illumination: selecting different neutral references can lead to different white-balance results.}
\label{fig:white-balance-mixed}
\end{figure}

\subsection{Von Kries Adaptation and a Gray Card}

The lecture presents a channel-gain model: multiply each channel by a separate factor. If a known neutral reference is recorded as $(r_w,g_w,b_w)$, choose gains proportional to
\begin{equation}
  g_R=\frac{1}{r_w}, \qquad g_G=\frac{1}{g_w}, \qquad g_B=\frac{1}{b_w}.
  \label{eq:von-kries}
\end{equation}
\mathlegend{
  \item $(r_w,g_w,b_w)$: measured RGB values of a white or gray reference object;
  \item $g_R,g_G,g_B$: per-channel multiplicative gains (the reference channels must be nonzero).
}
The slides also note that a light source can have a more complex effect than independent channel scaling; in general it may be represented by a full $3\times3$ transformation.

\subsection{White Balance without a Gray Card}

When no neutral reference is provided, the camera must infer the illuminant. The lecture lists several strategies:
\begin{itemize}
  \item \textbf{gray-world assumption}: assume the image averages $(r_{ave},g_{ave},b_{ave})$ should be gray and use inverse-average gains;
  \item \textbf{brightest-pixel assumption}: highlights often approximate the illuminant color, so use gains inversely proportional to the brightest channel values;
  \item \textbf{gamut mapping}: match the convex hull of observed image colors to a typical gamut under white light;
  \item \textbf{image statistics and learning}: estimate illuminant from learned regularities.
\end{itemize}

A higher-level method cited in the lecture performs \textbf{white balance by recognition}: for semantic classes detected in the image, estimate the illuminant that maps the class's measured average color toward the average color of that class in a database of typical images (Van de Weijer, Schmid, and Verbeek, ICCV 2007).

\FloatBarrier
\begin{chaptersummary}
Color constancy is a perceptual normalization problem; white balance is its computational analogue in imaging. Simple methods use channel gains and neutral references, while gray-world, brightest-pixel, gamut, statistical, and semantic-class methods estimate the illuminant when no reference card is available.
\end{chaptersummary}

\chapter{Color in Computer Vision and Texture}

\section{Uses of Color}

The lecture surveys several uses of color in computer vision and associates them with representative research examples:
\begin{table}[htbp]
\centering
\caption{Color-based computer-vision applications and references cited in the lecture.}
\label{tab:color-applications}
\begin{tabularx}{\textwidth}{@{}Y Y@{}}
\toprule
\rowcolor{TableTint}
\textbf{Application} & \textbf{Representative source named in the lecture} \\
\midrule
Color histograms for indexing and retrieval & Swain and Ballard, \emph{Color Indexing}, IJCV 1991 \\
Skin detection & Jones and Rehg, \emph{Statistical Color Models with Application to Skin Detection}, IJCV 2002 \\
Image segmentation and retrieval & Carson, Belongie, Greenspan, and Malik, \emph{Blobworld}, 1999 \\
Robot soccer & Sridharan and Stone, color calibration at RoboCup, 2006 \\
Appearance models for tracking & Ramanan, Forsyth, and Zisserman, \emph{Tracking People by Learning their Appearance}, PAMI 2007 \\
Judging visual realism & Lalonde and Efros, color compatibility for image realism, ICCV 2007 \\
\bottomrule
\end{tabularx}
\end{table}

\section{Structural and Statistical Texture}

The lecture defines \textbf{texture} as a description of the spatial arrangement of color or intensities in an image or selected region. A structural approach models a repeated arrangement of basic texture elements (texels). Natural images often make explicit texel extraction difficult or impossible, motivating statistical texture descriptions computed directly from gray tones or colors.

\begin{keyconcept}[Why statistical texture?]
Statistical texture descriptors are less intuitive than explicit texel models but are computationally efficient and can support both classification and segmentation. They are especially useful when natural texture elements cannot be segmented reliably.
\end{keyconcept}

\section{Local Binary Patterns}

For each pixel $p$, the Local Binary Pattern (LBP) example in the lecture compares the eight neighbors with the center value. For neighbor $i$,
\begin{equation}
  b_i=\begin{cases}
  0, & \text{if neighbor $i$ has value $\le p$},\\
  1, & \text{otherwise},
  \end{cases}
  \label{eq:lbp-bit}
\end{equation}
Here $p$ denotes the center intensity in this comparison, $i=1,\ldots,8$ follows the source neighbor ordering, and $b_i$ is the resulting comparison bit. Then the eight bits $b_1\ldots b_8$ form a local code. A region is represented by a histogram of these codes.

\begin{examplebox}{LBP lecture example}
For the $3\times3$ neighborhood
{\large
\[
\begin{bmatrix}
100&101&103\\
40&50&80\\
50&60&90
\end{bmatrix},
\]}
the center pixel is $50$. Starting at the upper-left neighbor and proceeding clockwise gives $100,101,103,80,90,60,50,40$. Comparing the eight neighbors with the center produces the bit pattern shown in the lecture, $11111100$, under its stated neighbor ordering. The histogram of such codes becomes a texture descriptor.
\end{examplebox}

The slides show LBP-histogram retrieval in the Fids image database and explicitly warn that low-level similarity does not always imply semantic similarity. Additional examples mentioned include deciduous-tree classification, finding cancer in Pap smears, and identifying regions of interest in breast and melanoma biopsy slides.

\section{Image Data Structure}

The lecture closes with a simple implementation-oriented image representation:\par
\begin{minipage}{\linewidth}
\begin{verbatim}
typedef struct{
    int h, w, c;
    float *data;
} image;
\end{verbatim}
\end{minipage}\par
Here $h$ is height, $w$ is width, and $c$ is the number of channels ($c=3$ for RGB and $c=1$ for grayscale in the example). The image data are stored as a linear array of floating-point values in $[0,1]$, i.e., a three-dimensional image tensor is flattened into one-dimensional memory.

\FloatBarrier
\begin{chaptersummary}
Color supports retrieval, segmentation, detection, tracking, calibration, and realism assessment. Texture describes spatial appearance; statistical descriptors such as LBP summarize local comparisons as histograms when explicit texel extraction is impractical.
\end{chaptersummary}

\par\vspace{1em}
\phantomsection
\addcontentsline{toc}{chapter}{Lecture 3 Quick Review}
\markboth{Lecture 3 Quick Review}{Lecture 3 Quick Review}
\begin{center}
{\sffamily\fontsize{22}{26}\selectfont\bfseries\color{Ink}Lecture 3 Quick Review\par}
\vspace{0.35em}{\color{SapienzaRed}\rule{0.82\textwidth}{0.8pt}}
\end{center}

\begin{multicols}{2}
\normalsize
\reviewgroup{Light and human vision}
\begin{enumerate}[label=\textcolor{SapienzaRed}{\bfseries\arabic*.},leftmargin=1.8em,itemsep=0.40em]
\item Color is a perceptual result of physical light interacting with the visual system.
\item Source spectra and surface reflectance jointly determine the spectral signal reaching the observer.
\item Rods are emphasized for intensity sensitivity; three cone classes support color vision.
\item Metamers have different spectra but the same perceived color under the relevant conditions.
\item Trichromacy allows color to be encoded by three primary weights.
\end{enumerate}
\reviewgroup{Color matching and spaces}
\begin{enumerate}[resume,label=\textcolor{SapienzaRed}{\bfseries\arabic*.},leftmargin=1.8em,itemsep=0.40em]
\item Grassmann's laws make additive color matching linear.
\item Color matching can be written compactly as $\mathbf e=C\mathbf t$.
\item RGB is a linear primary space; some spectral matches require subtractive components.
\item CIE XYZ uses nonnegative matching functions and chromaticity coordinates $x,y$.
\item McAdam ellipses show that equal $x,y$ distances are not equal perceptual differences.
\item HSV separates hue, saturation, and value for perceptually meaningful manipulation.
\end{enumerate}
\reviewgroup{Constancy and appearance}
\begin{enumerate}[resume,label=\textcolor{SapienzaRed}{\bfseries\arabic*.},leftmargin=1.8em,itemsep=0.40em]
\item Color/lightness constancy depends on context and adaptation.
\item White balance compensates for illuminant color; a gray card provides a direct neutral reference.
\item Gray-world, brightest-pixel, gamut, and learning methods estimate white balance without a reference card.
\item Texture describes spatial organization; statistical descriptors avoid explicit texel segmentation.
\item LBP compares neighbors with a center pixel and summarizes the resulting codes with a histogram.
\end{enumerate}
\end{multicols}


\end{document}
